\EMASCtitlepage{Motivação}

\begin{frame}[t]{Por que Ciência de Dados?}
  \vspace{3mm}
  A Ciência de Dados permite transformar informações em análises que ajudam a compreender problemas e tomar decisões.
  \medskip
  \begin{itemize}
    \item \textbf{Python:} organizar e processar os dados.
    \item \textbf{Estatística:} resumir, comparar e interpretar informações.
    \item \textbf{Visualização:} identificar padrões e comunicar resultados.
  \end{itemize}
  \begin{EMASCnote}
    Essas ferramentas são complementares. A partir de um problema real, podemos utilizar os dados disponíveis para investigar suas possíveis causas e propor soluções.
  \end{EMASCnote}
\end{frame}

\begin{frame}[t]{Ciência de Dados na prática: UFS e TRE}
  \vspace{2mm}
  \begin{columns}[T,onlytextwidth]
    \begin{column}{.65\textwidth}
      \textbf{Problema:} investigar quais urnas apresentavam atrasos nas Eleições Gerais de Sergipe.
      \medskip
      \begin{itemize}
        \item Organizar os registros das urnas eletrônicas.
        \item Analisar os tempos de votação e comparar as urnas.
        \item Visualizar os resultados em um painel interativo.
      \end{itemize}
      \begin{EMASCnote}
        Como identificar, a partir dos dados, quais urnas apresentaram os maiores atrasos?
      \end{EMASCnote}
    \end{column}
    \begin{column}{.32\textwidth}
      \centering
      {\bfseries\color{emascPurple} Dashboard UFS/TRE\par}
      \vspace{2mm}
      \includegraphics[width=31mm]{../motivacao/figuras/qr-code-dashboard-ufs-tre.png}\par
      \vspace{2mm}
      {\small\href{https://dashboard-ufs-tre.streamlit.app/}{\textcolor{emascPurple}{\textbf{Acessar painel}}}\par}
      \vspace{2mm}
      {\scriptsize\texttt{dashboard-ufs-tre\allowbreak.streamlit.app}\par}
    \end{column}
  \end{columns}
\end{frame}
